\section{Recovery-Aware Agent Control}

Recovery-Aware Agent Control (\RAAC) is a model-agnostic control layer for
tool-using agents. It observes interaction events, not benchmark condition
labels. The controller is forbidden from receiving the intervention family,
intervention identifier, hidden expected behaviour, answer key, evaluator
metadata, or gold. Its input is limited to observable signals such as tool
errors, timeout, malformed or partial output, missing required fields, schema
mismatch, contradictory observations, stale timestamps, inconsistent repeated
results, impossible values, insufficient evidence, unverifiable success, and
exhausted budgets.

\paragraph{State and action contract.}
The typed state machine moves through planning, acting, observation validation,
anomaly detection, bounded retry, alternate routing, cross-checking,
clarification, abstention, final verification, answering, and termination.
Illegal transitions fail closed. Every decision records its reason code,
trigger, remaining model/tool/token/time budgets, action, trace index, and
evidence class. A termination rule and monotone budget consumption prevent
infinite recovery loops.

\paragraph{Policies and fairness.}
\texttt{RAAC\_LIGHT} permits one verification, one retry, and at most one
alternate route for limited T4 compute. \texttt{RAAC\_FULL} permits additional
bounded contradiction resolution, clarification, abstention, and final
verification. Component ablations remove verification, retry, abstention,
cross-checking, alternate routing, or final verification. Equal-budget mode is
the primary causal comparison; practical-budget mode is a separate deployment
analysis. Both report clean and intervention success, extra model/tool calls,
tokens, latency, and measured cost.

\paragraph{Evidence status.}
Deterministic fixtures exercise clean success, transient and persistent failure,
conflict, stale evidence, malformed and partial output, premature success,
insufficient evidence, clarification, abstention, and alternate-route recovery.
They validate controller mechanics only. The hypothesis that \RAAC improves
robustness or preserves clean performance remains to be tested on locked,
human-validated tasks with audited non-oracle agents.
